\section{Indexed constructions: binding and annotated normalization}
\label{sec:annex-p07-ch88}

\begin{annexitem}{annex:p7.ch88.index-operations}{thm:p07-ch88-index-operations}{Scoped arithmetic in both binder fields}
\begin{proof}
Induct on the scoped grammar in both classes. Relocation fixes indices at most its cutoff and otherwise raises them by one, hence sends a legal depth-\(n\) index to one at depth \(n+1\). The star is fixed. Application uses the two object inductions. In either binder the annotation has the old cutoff and its body cutoff increased by one; their new depths are \(n+1\) and \(n+2\). The body class stays object or logical context as appropriate. Iteration gives \(\operatorname{add}(j,A)\) at depth \(m+j\) for \(A\) at depth \(m\). Lookup of \(A_k\), originally at depth \(n-k\), therefore has depth \(n\).

For substitution prove the stronger cutoff assertion: if \(N\in\mathcal A_i^n\) and \(Q\in\mathcal A_b^{n+1+i}\), then \(\operatorname{sub}_i(N,Q)\in\mathcal A_b^{n+i}\).
Protected indices remain legal; the selected index becomes \(\operatorname{add}(i,N)\), an object at depth \(n+i\); a greater index decreases by one and stays at most \(n+i\). The atom is fixed. Application keeps its two object children. Each binder uses the cutoff-\(i\) annotation induction and cutoff-\(i+1\) body induction; their output depths differ by one and their body class is preserved. Set \(i=0\). The object-argument hypothesis is essential at the selected variable.

Every recursive call visits a proper child of the original finite term. Relocation iteration decreases its finite counter; substitution does not recursively visit inserted images. Increasing a protected cutoff under a binder is not increasing the visited syntax. This proves totality without normalization. In the example, lookup raises \(1\) to \(2\); substitution lowers annotation \(3\) to \(2\), while the body at cutoff one replaces selected \(2\) by \(\operatorname{add}(1,1)=2\).
\end{proof}
\end{annexitem}

\begin{annexitem}{annex:p7.ch88.binding-transport}{thm:p07-ch88-binding-transport}{The actual binding and substitution encoding}
\begin{proof}
For an ordered external variable vector \(\vec z=(z_1,\ldots,z_n)\)
from oldest to newest, define \(T_{\vec z}(\ast)=\ast\) and
\(T_{\vec z}(k)=z_{n+1-k}\). Map application homomorphically.
Map a source product to \(\Pi z:T_{\vec z}(A).
T_{\vec z,z}(B)\), and a lambda to the analogous annotated lambda,
choosing \(z\) fresh for the external vector. An environment
\(\Gamma=[x_n:A_n]\cdots[x_1:A_1]\ast\) maps to the ordered
declarations whose annotations are translated in their respective
older prefixes, rather than to a single type term. A logical context
appearing as a type \(\Delta\) still maps by \(T\) to a product
ending in \(\ast\). These are distinct roles of the source's
same product syntax and are not conflated.

The named encoding has a unique indexed inverse on its image, modulo
alpha-renaming: each variable names its position in the external vector
or binder stack; products/lambdas extend that stack only for the body,
and application recursively decodes both children. Induction on syntax
proves both composites are identity on the scoped source grammar and
its named image. It is not claimed to be onto arbitrary raw CC terms,
which additionally contain \(\Box\) and many ill-typed forms.

Corrected relocation commutes with this encoding. More generally, raise
the external indices by \(k\) with the correct cutoff under every
internal binder. A variable then still designates the same older named
variable; an internal bound reference still designates its binder.
The star/application cases are immediate; annotation/body induction
proves the binder cases. In particular the source lookup
\(\Gamma/l=\operatorname{add}(l,A_l)\) translates to the
original named annotation of the corresponding declaration, not to
an annotation containing that declaration's own fresh name.

For source \(Q\) scoped over \(\Gamma[x:A]\) and \(N\) scoped
over \(\Gamma\), corrected substitution/removal of index one gives
\[
 T_{\vec z}([N/x]Q)
   =T_{\vec z,z}(Q)[T_{\vec z}(N)/z].
\]
For index one this gives the argument; a larger external index decreases
by one when the removed binder disappears, identifying the same external
name on both sides. Under an inner binder, its local references stay
fixed, the substituted argument is correctly relocated, and a fresh
named binder prevents capture. Induction on annotation and body proves
both product and lambda cases, and application combines the child
identities. The same structural proof gives simultaneous substitutions
in all later environment annotations. A source beta contraction thus
maps to one designated named beta contraction at the same position,
including a contraction in a binder annotation or product body.
\end{proof}
\end{annexitem}

\begin{annexitem}{annex:p7.ch88.full-transport}{thm:p07-ch88-full-transport}{Joint validity typing and equality with sorted conversion}
\begin{proof}
Use the structural CC facts proved with \cref{thm:p12-har-cc-normalization} in Annex~\ref{annex:p12.har.cc-normalization}: validity/projection, generation/type correctness, context conversion, substitution, uniqueness, raw beta confluence and compatible beta subject reduction. Those facts do not assume normalization. A typable subject cannot beta-convert to the normal atom \(\Box\), since confluence and subject reduction would give that atom a forbidden typing. Source images have no literal \(\Box\). We work throughout with the stated corrected binding convention and actual full-source rule replacements.

Write \(\widehat\Gamma\) for the named environment just defined.
Simultaneously induct on original full-system derivations to establish:
\(\Gamma\vdash\Delta\) gives
\(\widehat\Gamma\vdash_{\rm CC}T\Delta:\Box\);
\(\Gamma\vdash M:P\) gives
\(\widehat\Gamma\vdash_{\rm CC}TM:TP\);
and \(\Gamma\vdash M\cong N\) gives raw
\(TM=_\beta TN\), together with the following endpoint transport:
for every target type \(U\), either endpoint's CC judgment at \(U\)
implies the other endpoint's judgment at that same \(U\).
This stronger equality induction supplies the target sort premise needed
for source conversion; it does not assume the source's Lemmas~3--4,
whose full proofs the article refers elsewhere.

The primitive validity axiom maps to \(\ast:\Box\) in the empty
target environment. Extension by a source logical context \(\Delta\)
uses its mapped \(\Box\) judgment as the annotation certificate;
extension by an object \(M:\ast\) uses its mapped \(\ast\)
judgment. Both make the extended target environment valid, and the
weakened axiom types \(\ast\) there. A validity product rule has
premise \(\Gamma[x:A]\vdash\Delta\); its induction gives
\(T\Delta:\Box\) in the extended target context. CC validity
extracts the annotation certificate \(TA:s_1\); the triple
\((s_1,\Box,\Box)\) forms its image product. An object product rule
has body premise \(M:\ast\); similarly use
\((s_1,\ast,\ast)\). Both possible domain sorts are admitted.

For source projection, its validity premise makes
\(\widehat\Gamma\) valid. CC projection types the corresponding
named variable at its older-prefix annotation, and \cref{thm:p07-ch88-binding-transport}'s lookup
identity identifies this with \(T(\Gamma/l)\). A source abstraction
uses its body induction in the target extension. Its annotation sort
is obtained from that extension's CC validity. Its body type \(TP\)
has a sort by CC type correctness; if it is the literal sort \(\ast\),
use \(\ast:\Box\). It cannot be \(\Box\), which has no source
preimage. The requisite product is therefore well formed by one of
the four triples, and CC abstraction applies. Application uses the
two mapped premises and \cref{thm:p07-ch88-binding-transport}'s substitution identity.

Full equality reflexivity (for typed objects or valid logical contexts)
and symmetry/transitivity map to raw beta equality and preserve endpoint
typing immediately or by composing the endpoint inductions. For product
congruence, suppose an endpoint product is typed at \(U\). CC generation
exposes its annotation sort and body sort. Annotation equality induction
transports the first sort judgment to the second annotation; body
equality induction transports the body sort in the old binder context.
CC context conversion moves it to the new, beta-equal annotation.
Rebuild the product at its original result sort, then restore \(U\)
by conversion when required. If \(U=\Box\), the generated product
result is literally \(\Box\), so no illegal top-sort conversion is
used. For the reverse direction, first transport the new domain sort
to the old annotation, then use context conversion to carry the
second body judgment into the old binder context, where the body
equality induction is stated. Apply its reverse implication there
and rebuild the first product. This preserves the required body
context rather than assuming the equality was supplied in both contexts.

For lambda congruence, generation of a typed endpoint gives its literal
body judgment and a sorted product type. Annotation equality transports
its domain sort; body equality transports its body at the same generated
body type. Context conversion transports that body judgment and its
body-type sort to the changed annotation. Abstraction produces the new
product type, beta-equal to the old one; convert back to \(U\).
No lambda has the unsorted top as a type, by the structural CC generation facts. For full application
congruence, endpoint generation gives function and argument judgments.
Their endpoint inductions retain the same function product and argument
domain. Apply them; raw argument equality and substitution give equality
of the new and old codomain instances, permitting conversion to \(U\).
Both directions use the same steps with endpoints interchanged.
This covers the full system's simultaneous function/argument replacement,
not only the restricted one-child congruence rules.
For reverse lambda transport, analogously carry the second endpoint's
generated body judgment and body-type sort from the new domain back to
the old domain before applying the reverse body equality implication.
Then reconstruct the first lambda and restore the given type. Compatible
raw beta equality of each mapped domain/body or function/argument pair
gives the asserted raw lambda, product, or application equality in
these congruence cases.

For the full beta equality generator, its source body/argument premises
are mapped by the joint induction. The same type-correctness/product
construction as in abstraction gives a CC-typed root redex. Substitution
types its contractum at the same codomain instance. Raw equality is the
mapped root beta step. For endpoint transport from any additionally
given root type \(U\), use CC subject reduction. In the reverse
direction, uniqueness between the contractum's principal instance and
its given type \(U\) gives beta equality; convert the already typed
root to \(U\). Here \(U\) is sorted: if it were \(\Box\), its
equality to the explicitly sorted principal result type would contradict
the structural CC result's impossibility of a typable subject converting to that atom.
Thus this case does not assume unsafe general beta expansion.

Finally consider the original typing conversion. Its subject induction
gives \(TM:TP\). CC type correctness sorts \(TP\), with the same
literal-star exception handled by \(\ast:\Box\); no source type
maps to \(\Box\). The equality endpoint induction applied to this
sort judgment gives \(TQ\) at the same sort, while its raw equality
gives \(TP=_\beta TQ\). CC conversion therefore licenses the desired
\(TM:TQ\). This closes the mutual induction, including the otherwise
missing target-type sort premise. No normalization entered transport.
\end{proof}
\end{annexitem}

\begin{annexitem}{annex:p7.ch88.restricted-inclusion}{thm:p07-ch88-restricted-inclusion}{All corrected restricted rules as full derivations}
\begin{proof}
First establish generation inside the full source surface. A derivation of \(\Gamma\vdash(MN):P\) has premises \(\Gamma\vdash M:[x:A]B\) and \(\Gamma\vdash N:A\) for some \(A,B\). Induct on that derivation: application supplies them; a final typing conversion has the unchanged subject and environment, so use its premise induction. Projection, abstraction and product formation have different subject constructors. There is no independent source weakening rule. Thus these premises are in the same environment even if the result was converted.

Induct jointly on corrected restricted derivations. Common validity/typing and typing conversion are unchanged full rules after the premise inductions. Context reflexivity, symmetry, transitivity and product congruence are unchanged. Restricted object reflexivity is full reflexivity with type \(\Delta\); logical lambda congruence is a full instance with \(N_1=\Delta\); logical beta is a full instance with \(P=\Delta\). Their logical hypotheses are retained.

For function-only compatibility, generation of the mapped full application typing supplies the argument typing, hence full reflexivity \(N\cong N\). Combine that, the mapped function equality and the application typing in the simultaneous full rule. For argument-only compatibility, generation supplies the function typing, hence \(M\cong M\); the same full rule concludes the required argument-only equality. These are all restricted cases. The construction preserves the indexed syntax and context, not a converse equality of the systems.
\end{proof}
\end{annexitem}

\begin{annexitem}{annex:p7.ch88.logical-transport}{thm:p07-ch88-logical-transport}{Primitive logical clauses and positive transitive paths}
\begin{proof}
Induct on the directed derivation. The logical beta generator maps to the one root contraction of \cref{thm:p07-ch88-binding-transport}; its guards introduce no new rewrite. Product and lambda body clauses lift the positive body induction through the corresponding fresh named binder. Object-product, logical-product and lambda annotation clauses lift their positive annotation induction through that exact field with the body unchanged. The body's fresh name is the same for old and new annotations. Application clauses lift through their respective left or right child. The original typing guards stay in the source derivation and are not promoted to zero-step target edges.

Transitivity concatenates its two induction paths, both nonempty, so their concatenation is positive. These exhaust the directed sheet. No normalization is used and a transitive derived edge is not asserted to be one contraction.
\end{proof}
\end{annexitem}

\begin{annexitem}{annex:p7.ch88.full-normalization}{thm:p07-ch88-full-normalization}{Full annotated open normalization from proved CC}
\begin{proof}
Full transport gives a CC-typed named image for an object typing or a kind judgment at \(\Box\) for a valid logical-context term. The independently proved exact CC theorem \cref{thm:p12-har-cc-normalization} applies to both in their valid open contexts. If an infinite compatible annotated source beta path existed, scoped substitution would keep every contractum in its source class and external depth. Use the same external vector throughout. Binding transport sends every root, application-child or binder annotation/body contraction to one target contraction. This yields an infinite target path, contradicting the CC theorem.

No closed inhabitant of every annotation is assumed: the proved target theorem includes valid open contexts via identity raw valuations and canonical semantic values. No source field or type argument was erased. This is strong normalization of all compatible paths, not just a normalizing strategy or weak-head computation.

A corrected restricted validity/typing first has the full derivation of \cref{thm:p07-ch88-restricted-inclusion}; the same argument applies. For an infinite directed logical chain, its initial image is typed as above and every edge contributes a positive finite beta path by \cref{thm:p07-ch88-logical-transport}. Concatenating those paths contradicts the target theorem again. No source checking correctness or source normalization premise was used; the original closed stripped-image corollary remains separately cited and is not the annotation-termination argument.
\end{proof}
\end{annexitem}
